\section{Finite unit registers and protected coordinates}
\label{a-finite-units}

Adic refinement first changes the address of a state. Its amplitudes remain
complex numbers, and its measurement probabilities remain ordinary Born
probabilities. For a finite address set $X$ we use the counting inner product
on $\ell^2(X;\C)$, so the address vectors $|x\rangle$ are orthonormal.
A bijection of address sets therefore gives a permutation unitary without
changing the quantum probability rule. Arithmetic determines the address
factorization. A Hamiltonian and a measurement prescription supply the
additional physical choices.

The first coordinate issue is specific to units. For an odd prime $p$ and
an integer $h\geq0$, set
\[
 U_{p,h}=(\Z/p^{h+1}\Z)^\times.
\]
For $p=2$ the useful indexing starts one digit later:
\[
 U_{2,h}=(\Z/2^{h+2}\Z)^\times.
\]
The shifted dyadic index retains the sign even when the principal phase
register has dimension one. It prevents the two units modulo four from
being collapsed into one phase coordinate.

\subsection{Torsion and principal units}

\begin{proposition}[Finite unit coordinates]
\label{a-finite-unit-coordinates}
For odd $p$, choose a unit $\zeta_{p,h}$ of order $p-1$ whose reduction
generates $(\Z/p\Z)^\times$. The maps
\begin{align}
 (t,a)&\longmapsto\zeta_{p,h}^{\,t}(1+p)^a
       \pmod{p^{h+1}},
 & (t,a)&\in\Z/(p-1)\Z\times\Z/p^h\Z,
 \label{a-finite-odd-coordinate}\\
 (s,a)&\longmapsto(-1)^s5^a\pmod{2^{h+2}},
 & (s,a)&\in\Z/2\Z\times\Z/2^h\Z
 \label{a-finite-dyadic-coordinate}
\end{align}
are group isomorphisms onto $U_{p,h}$ and $U_{2,h}$ respectively.
\end{proposition}

\begin{proof}
For odd $p$ the valuation identity
\[
 v_p\bigl((1+p)^{p^k}-1\bigr)=k+1
 \qquad(k\geq0)
\]
shows that $1+p$ has order $p^h$ modulo $p^{h+1}$. To check the
identity, the case $k=0$ is immediate. If $u=1+c p^{k+1}$ with
$p\nmid c$, the binomial expansion of $u^p-1$ has first term
$c p^{k+2}$. Every remaining term has strictly larger $p$-valuation
when $p$ is odd. This proves the induction. The subgroup generated by
$1+p$ is contained in the kernel of reduction modulo $p$ and has the
same cardinality $p^h$, so it is that entire kernel.

For completeness, an order-$p-1$ lift of any nonzero residue can be chosen
coherently. The root equation $x^{p-1}=1$ has nonzero derivative
$(p-1)x^{p-2}$ modulo $p$. If $x_n$ is a root modulo $p^n$, substituting
$x_n+c p^n$ modulo $p^{n+1}$ determines one and only one residue $c$
that cancels the error. Starting with a generator of
$(\Z/p\Z)^\times$ produces $\zeta_{p,h}$ with the required order.
The cyclicity of this residue group follows from the elementary fact that
a degree-$d$ polynomial over a field has at most $d$ roots: in a finite
subgroup of the multiplicative group, the least common multiple $m$ of
the element orders is realized by an element. To see this, take for each
prime divisor of $m$ a component having the maximal prime-power order
and multiply those commuting components of coprime orders. All elements solve
$x^m=1$. Consequently the subgroup has size at most $m$, and Lagrange's
theorem gives the reverse bound. The subgroup is cyclic.

If $\zeta_{p,h}^{\,t}(1+p)^a=1$, reduction modulo $p$ gives
$t=0$ in $\Z/(p-1)\Z$. The principal-unit order then gives
$a=0$ in $\Z/p^h\Z$. Thus the displayed map is injective, and its
domain has size $(p-1)p^h=|U_{p,h}|$, which proves surjectivity.

For the dyadic case, successive squaring gives
\[
 v_2\bigl(5^{2^k}-1\bigr)=k+2.
\]
Indeed $5-1$ has valuation two, and for every such power the factor
$5^{2^k}+1$ has valuation one. Thus $5$ has order $2^h$ modulo
$2^{h+2}$. Its powers are exactly the $2^h$ units congruent to one
modulo four. Multiplication by $-1$ supplies the units congruent to
three modulo four. These two classes are disjoint, giving the claimed
unique coordinates.
\end{proof}

In the odd-prime register $t$ is a torsion coordinate and $a$ is a
principal phase coordinate. In the dyadic register $s$ is a sign bit and
$a$ is again the principal phase. The Hilbert-space factorizations are
therefore
\begin{align}
 \ell^2(U_{p,h})&\cong
 \C^{p-1}_{\mathrm{tors}}\otimes\ell^2(\Z/p^h\Z),
 &&p\text{ odd},\nonumber\\
 \ell^2(U_{2,h})&\cong
 \C^2_{\mathrm{sign}}\otimes\ell^2(\Z/2^h\Z).
 \label{a-finite-unit-hilbert}
\end{align}
These are unitary changes of basis. They do not assert that a generic
state is a product state, or that a generic Hamiltonian is a sum of
independent register Hamiltonians. Entanglement and interactions are
preserved under the change of coordinates.

The lifting construction in the proof permits the coherent choice
$\zeta_{p,h+1}\equiv\zeta_{p,h}\pmod{p^{h+1}}$. With that choice,
reduction of unit addresses retains $t$ and reduces $a$ modulo $p^h$.
If the torsion lifts are instead chosen independently at each depth, the
finite factorizations remain exact, but a comparison between levels must
include the resulting coordinate changes. A levelwise permutation is
not by itself a specified refinement tower.



\subsection{CRT readdressing and the Fourier dual}

Let $\mathcal P_{\mathrm o}$ be a finite set of odd primes and put
\[
 M=2^{h_2+2}\prod_{p\in\mathcal P_{\mathrm o}}p^{h_p+1}.
\]
The Chinese remainder theorem gives a bijection from the unit addresses
modulo $M$ to their local unit addresses. Combining it with
Proposition~\ref{a-finite-unit-coordinates} gives a permutation unitary
\begin{equation}
 \mathcal U_{\mathrm{CRT}}:
 \ell^2\bigl((\Z/M\Z)^\times\bigr)
 \longrightarrow
 \C^2_{\mathrm{sign}}\otimes\ell^2(\Z/2^{h_2}\Z)
 \otimes\bigotimes_{p\in\mathcal P_{\mathrm o}}
 \left(\C^{p-1}_{\mathrm{tors}}\otimes\ell^2(\Z/p^{h_p}\Z)\right).
 \label{a-finite-crt}
\end{equation}
For each basis vector, $\mathcal U_{\mathrm{CRT}}|u\rangle$ is the basis
vector with its unique tuple of coordinates. Distinct units have distinct
tuples, so inner products are preserved. The dimensions also agree:
\[
 \varphi(M)=2^{h_2+1}\prod_{p\in\mathcal P_{\mathrm o}}(p-1)p^{h_p}.
\]
Multiplication of units becomes addition in each cyclic coordinate.

We use the positive-phase Fourier convention
\begin{equation}
 (F_n)_{k x}=n^{-1/2}\exp(2\pi\mathrm{i}kx/n).
 \label{a-finite-fourier-convention}
\end{equation}
The character transform of the unit group consequently factors as
\begin{equation}
 F_{U_M}=\mathcal U_{\mathrm{CRT}}^*
 \left[F_2\otimes F_{2^{h_2}}\otimes
 \bigotimes_{p\in\mathcal P_{\mathrm o}}(F_{p-1}\otimes F_{p^{h_p}})
 \right]\mathcal U_{\mathrm{CRT}}.
 \label{a-finite-unit-fourier}
\end{equation}
Here the unit-group character labels have been identified with the same
cyclic coordinate sets. This is a transform on the multiplicative unit
group, rather than the additive Fourier transform on all residues
modulo $M$.

The additive transform has a related factorization, with different
permutations on its input and its dual output. Let $a,b$ be positive
coprime integers and $N=ab$. Define
\begin{align*}
 \mathcal P_{\mathrm{in}}|n\rangle
 &=|n\bmod a\rangle\otimes|n\bmod b\rangle,\\
 \mathcal P_{\mathrm{out}}|k\rangle
 &=|k b^{-1}\bmod a\rangle\otimes|k a^{-1}\bmod b\rangle,
\end{align*}
where the inverses are taken in the indicated residue rings. Then
\begin{equation}
 F_N=\mathcal P_{\mathrm{out}}^*(F_a\otimes F_b)
       \mathcal P_{\mathrm{in}}.
 \label{a-finite-additive-fourier}
\end{equation}
To verify the phase, let $n_a=n\bmod a$ and $n_b=n\bmod b$. At the
dual coordinates $k_a=kb^{-1}\bmod a$ and $k_b=ka^{-1}\bmod b$, the
product kernel is
\[
 N^{-1/2}\exp\!\left[2\pi\mathrm{i}
      \left(\frac{k_a n_a}{a}+\frac{k_b n_b}{b}\right)\right].
\]
The numerator $b k_a+a k_b$ is congruent to $k$ modulo both $a$ and
$b$, hence modulo $ab$. This kernel equals
$N^{-1/2}\exp(2\pi\mathrm{i}kn/N)$, proving
\eqref{a-finite-additive-fourier}. Omitting the dual-output permutation
would generally give a different phase labeling.

\subsection{What is lost when the sign is removed}

The linear address-forgetting map
\begin{equation}
 C_{\mathrm{phase}}|s,a\rangle=|a\rangle
 \label{a-finite-forgetting}
\end{equation}
maps a $2^{h+1}$-dimensional dyadic unit register onto a
$2^h$-dimensional principal phase register. Its kernel is the span of
the $2^h$ vectors $|0,a\rangle-|1,a\rangle$. Thus
\[
 \operatorname{rank}C_{\mathrm{phase}}=2^h,
 \qquad\dim\ker C_{\mathrm{phase}}=2^h.
\]
This is a rank-deficient linear sum of amplitudes. It is not an
isometry, and it is not the quantum channel that discards a subsystem.
For example it sends the normalized state
$(|0,a\rangle+|1,a\rangle)/\sqrt2$ to $\sqrt2|a\rangle$.
Its adjoint product therefore does not equal the identity.

Discarding the sign as a quantum operation instead takes the partial
trace $\rho\mapsto\operatorname{Tr}_{\mathrm{sign}}\rho$. Both pure
states $|0,a\rangle$ and $|1,a\rangle$ then produce the same reduced
phase state $|a\rangle\langle a|$. Yet the sign observable
\[
 Z_{\mathrm{sign}}=(|0\rangle\langle0|-|1\rangle\langle1|)\otimes I
\]
has expectations $+1$ and $-1$ on those states. No measurement on the
retained phase state can reconstruct this difference. The obstruction is
information loss in a stated map, not an alternative probability rule.

A separate arithmetic example shows why integral residue selectors and
rational spectral idempotents should be distinguished. Let
\[
 A=\begin{pmatrix}1&1\\0&-1\end{pmatrix},\qquad A^2=I,
 \qquad E_\pm=\frac{I\pm A}{2}.
\]
The identities $E_\pm^2=E_\pm$, $E_+E_-=0$, and $E_++E_-=I$
follow directly from $A^2=I$. Both idempotents contain a denominator
two, and the resultant of $x-1$ and $x+1$ has dyadic valuation one.
Modulo two the two roots agree and $(A-I)^2=0$, while $A-I$ is
nonzero. Thus the residual matrix can retain a nilpotent extension even
though it no longer distinguishes the two rational eigenvalues.
These $E_\pm$ are complementary algebraic spectral idempotents. Since
$A$ is not self-adjoint in the standard counting inner product, they
are not orthogonal quantum measurement projectors in that inner product.
The unit-address sign projector, by contrast, is an ordinary diagonal
orthogonal projector.



\section{Cylinder averages, Haar details, and digit analysis}
\label{a-finite-haar}

The principal phase register has a canonical finite filtration. For a
prime $p$ and depth $h\geq0$, set
\[
 G_{p,h}=\Z/p^h\Z,\qquad\mathcal H_{p,h}=\ell^2(G_{p,h};\C).
\]
At depth zero this is a one-dimensional space. A level-$r$ cylinder
contains all residues with a prescribed prefix modulo $p^r$, and has
$p^{h-r}$ members. Throughout this section depth refers to the number of
retained least significant base-$p$ digits.

\subsection{Orthogonal averaging and its detail resolution}

For $0\leq r\leq h$, define the cylinder average
\begin{equation}
 (E_r^{(h)}f)(x)=p^{-(h-r)}
       \sum_{y\equiv x\, (\mathrm{mod}\,p^r)}f(y).
 \label{a-finite-cylinder-average}
\end{equation}
Its matrix is real and symmetric. Averaging twice over the same cylinder
has the same effect as averaging once. Thus $E_r^{(h)}$ is the
orthogonal projector onto functions constant on level-$r$ cylinders.
The range has dimension $p^r$.

\Needspace{16\baselineskip}
\begin{proposition}[Finite Haar resolution]
\label{a-finite-haar-resolution}
The cylinder projectors satisfy
\[
 E_r^{(h)}E_s^{(h)}=E_{\min(r,s)}^{(h)},\qquad E_h^{(h)}=I.
\]
Define
\[
 D_0^{(h)}=E_0^{(h)},\qquad
 D_r^{(h)}=E_r^{(h)}-E_{r-1}^{(h)}\quad(1\leq r\leq h).
\]
Then the $D_r^{(h)}$ are mutually orthogonal projectors and
\begin{align}
 \sum_{r=0}^hD_r^{(h)}&=I,
 &D_r^{(h)}D_s^{(h)}&=\delta_{rs}D_r^{(h)},\nonumber\\
 \operatorname{rank}D_0^{(h)}&=1,
 &\operatorname{rank}D_r^{(h)}&=(p-1)p^{r-1}\quad(r\geq1).
 \label{a-finite-detail-ranks}
\end{align}
\end{proposition}

\begin{proof}
If $r\leq s$, the level-$s$ cylinders partition each level-$r$
cylinder into equally sized groups. Averaging over the smaller groups
and then over the whole group returns the full average, giving
$E_rE_s=E_r$. The self-adjointness of both sides also gives
$E_sE_r=E_r$. This proves the minimum rule.

For $r\geq1$,
$(E_r-E_{r-1})^2=E_r-2E_{r-1}+E_{r-1}=D_r$.
For $1\leq r<s$, expansion of
$(E_r-E_{r-1})(E_s-E_{s-1})$ gives zero by the same minimum
rule. The products with $D_0$ vanish as well. Telescoping gives the
resolution of the identity. Finally $\Ran E_{r-1}\subset\Ran E_r$,
so the detail rank is $p^r-p^{r-1}$, while $E_0$ has rank one.
\end{proof}

The decomposition is exact at each finite depth:
\begin{equation}
 \mathcal H_{p,h}=\bigoplus_{r=0}^h\Ran D_r^{(h)},\qquad
 \norm f^2=\sum_{r=0}^h\norm{D_r^{(h)}f}^2.
 \label{a-finite-parseval}
\end{equation}
A depth cutoff $m\leq h$ retains $E_m^{(h)}=\sum_{r=0}^mD_r^{(h)}$.
Its discarded norm is exactly
\[
 \norm{(I-E_m^{(h)})f}^2=
 \sum_{r=m+1}^h\norm{D_r^{(h)}f}^2.
\]
No smoothness assumption is needed for this finite Parseval identity.
A rate of decrease with $m$ requires an additional bound on the detail
weights of the state.

\subsection{The child characters form the Haar basis}

Write $\omega_p=\exp(2\pi\mathrm{i}/p)$. For $1\leq r\leq h$,
parent prefix $a\in\{0,\ldots,p^{r-1}-1\}$, and character
$q\in\{1,\ldots,p-1\}$, define
\begin{equation}
 \psi_{r,a,q}^{(h)}(x)=
 \begin{cases}
 p^{-(h-r+1)/2}\omega_p^{q d_r(x)},
       &x\equiv a\pmod{p^{r-1}},\\
 0,&\text{otherwise},
 \end{cases}
 \label{a-finite-wavelet}
\end{equation}
where $d_r(x)$ is the digit multiplying $p^{r-1}$. Include the
constant scaling vector $\psi_0^{(h)}(x)=p^{-h/2}$.

Each wavelet is constant on level-$r$ cylinders and has zero average
on its parent level-$r-1$ cylinder, because
$\sum_{d=0}^{p-1}\omega_p^{qd}=0$ for nonzero $q$. Thus it lies in
$\Ran D_r^{(h)}$. Its support contains $p^{h-r+1}$ points, giving
unit norm. Distinct parents have disjoint supports. For a fixed parent
the inner product of the characters $q,q'$ is
\[
 p^{-1}\sum_{d=0}^{p-1}\omega_p^{(q'-q)d}=\delta_{qq'}.
\]
Wavelets at different levels are orthogonal because their detail
projectors are orthogonal. There are exactly
$(p-1)p^{r-1}$ wavelets at level $r$, equal to the rank in
\eqref{a-finite-detail-ranks}, so they span that detail subspace. Hence
\begin{equation}
 D_r^{(h)}=\sum_{a=0}^{p^{r-1}-1}\sum_{q=1}^{p-1}
       |\psi_{r,a,q}^{(h)}\rangle\langle\psi_{r,a,q}^{(h)}|.
 \label{a-finite-wavelet-projector}
\end{equation}
The omitted character $q=0$ is the parent average. For $p=2$ there is
one average and one difference per parent. For odd $p$ there is one
average and $p-1$ nontrivial child characters. The dyadic sign factor of
\eqref{a-finite-unit-hilbert} is independent of this average/detail
split in the principal phase register.

\adicfigure{02_haar_channels}{For $p=3$, the normalized child characters
contain one average and two orthogonal complex details. Solid curves
show real parts and dashed curves show imaginary parts. The level
rank $(p-1)p^{r-1}$ grows by the branching factor. The atlas uses
$\mathsf F_r$ for the corresponding infinite cylinder detail; its
finite restriction is $D_r^{(h)}$ for $r\leq h$.}



\subsection{The refinement isometry and its adjoint}

One more digit supplies $p$ children for every address. The uniform
refinement map is
\begin{equation}
 (V_{p,h}f)(a+d p^h)=p^{-1/2}f(a),
 \quad 0\leq a<p^h,\quad0\leq d<p.
 \label{a-finite-refinement}
\end{equation}
In tensor coordinates it is $f\mapsto|+_p\rangle\otimes f$,
where $|+_p\rangle=p^{-1/2}\sum_d|d\rangle$ and the child digit is
listed first. The norm calculation is
\[
 \norm{V_{p,h}f}^2=\sum_a\sum_d p^{-1}|f(a)|^2=\norm f^2.
\]
The adjoint is the normalized amplitude sum
\begin{equation}
 (V_{p,h}^*g)(a)=p^{-1/2}\sum_{d=0}^{p-1}g(a+d p^h).
 \label{a-finite-refinement-adjoint}
\end{equation}
Consequently
\begin{equation}
 V_{p,h}^*V_{p,h}=I,\qquad
 V_{p,h}V_{p,h}^*=E_h^{(h+1)},\qquad
 I-V_{p,h}V_{p,h}^*=D_{h+1}^{(h+1)}.
 \label{a-finite-refinement-projectors}
\end{equation}
The first composition is an identity on the coarse space. The second
is a projector on the fine space, not a fine-space identity.

Replicated vectors are constant in the new digit. All old cylinder
averages commute with replication, giving
\begin{align}
 E_r^{(h+1)}V_{p,h}&=V_{p,h}E_r^{(h)}\quad(0\leq r\leq h),
 \nonumber\\
 D_r^{(h+1)}V_{p,h}&=V_{p,h}D_r^{(h)}\quad(0\leq r\leq h),
 &D_{h+1}^{(h+1)}V_{p,h}&=0.
 \label{a-finite-detail-intertwining}
\end{align}
Repeated refinements compose as isometries. For $w\geq0$ and one
coarse address,
\begin{equation}
 V_{p,h+w-1}\cdots V_{p,h}|a\rangle
 =p^{-w/2}\sum_{\mathbf d\in\{0,\ldots,p-1\}^w}|a,\mathbf d\rangle,
 \label{a-finite-nested-refinement}
\end{equation}
with the empty product interpreted as the identity. The factor
$p^{-w/2}$ is forced by the isometry and orthonormality of the
descendant addresses.

\subsection{Digit reversal and the polyphase split}
\label{a-finite-polyphase}

For $h\geq1$, put $N=p^h$ and $m=p^{h-1}$. Two permutation
unitaries separate different digits:
\begin{align}
 \mathcal A^{\mathrm{std}}_{p,h}|pk+r\rangle
 &=|r\rangle\otimes|k\rangle,
 &0\leq r<p,\quad0\leq k<m,\nonumber\\
 \mathcal A^{\mathrm{ad}}_{p,h}|a+d p^{h-1}\rangle
 &=|d\rangle\otimes|a\rangle,
 &0\leq d<p,\quad0\leq a<m.
 \label{a-finite-digit-splits}
\end{align}
The standard polyphase split extracts the least significant residue.
The adic split extracts the newest, most significant digit at the
current depth. Let $\mathcal R_{p,h}$ reverse the $h$ base-$p$ digits,
with $\mathcal R_{p,0}=I$ on the one-dimensional space. Then
\begin{equation}
 \mathcal A^{\mathrm{ad}}_{p,h}=
 (I_p\otimes\mathcal R_{p,h-1})
 \mathcal A^{\mathrm{std}}_{p,h}\mathcal R_{p,h}.
 \label{a-finite-digit-reversal}
\end{equation}
Indeed reversing $d_0+d_1p+\cdots+d_{h-1}p^{h-1}$ makes
$d_{h-1}$ the least significant digit. Standard analysis removes that
digit, and reversing the remaining $h-1$ digits restores the prefix
$a=d_0+\cdots+d_{h-2}p^{h-2}$. This proves the identity on an
orthonormal basis.

Let the new-digit modulation be
$M_{p,h}|a+d p^{h-1}\rangle=\omega_p^d|a+d p^{h-1}\rangle$.
The selector of child digit $d$ is the exact roots-of-unity filter
\begin{equation}
 R_d=\frac1p\sum_{\ell=0}^{p-1}\omega_p^{-d\ell}M_{p,h}^{\ell}.
 \label{a-finite-root-selector}
\end{equation}
On a child with digit $e$, the coefficient is
$p^{-1}\sum_\ell\omega_p^{(e-d)\ell}=\delta_{de}$. Therefore
$R_dR_e=\delta_{de}R_d$, $R_d^*=R_d$, and $\sum_dR_d=I$.
These projectors select address digits. They are different from the
projectors onto Fourier characters of a child register.

Apply the child Fourier transform after adic analysis:
\begin{equation}
 \mathcal W_{p,h}=(F_p\otimes I_m)\mathcal A^{\mathrm{ad}}_{p,h}.
 \label{a-finite-child-fourier}
\end{equation}
The zero child-frequency row of $F_p$ is $p^{-1/2}(1,\ldots,1)$.
Hence
\begin{equation}
\begin{aligned}
 (\langle0|\otimes I_m)\mathcal W_{p,h}&=V_{p,h-1}^*,\\
 \mathcal W_{p,h}^*
 \bigl[(I_p-|0\rangle\langle0|)\otimes I_m\bigr]
 \mathcal W_{p,h}&=D_h^{(h)}.
\end{aligned}
\label{a-finite-child-detail}
\end{equation}
The second identity subtracts the coarse zero-frequency projector from
the identity using \eqref{a-finite-refinement-projectors}. Recursive
application to the zero-frequency channel alone constructs a unitary
Haar analysis map $\mathcal T_{p,h}$ whose channel projectors $\Pi_r$
have the ranks in \eqref{a-finite-detail-ranks} and obey
\begin{equation}
 \mathcal T_{p,h}^*\Pi_r\mathcal T_{p,h}=D_r^{(h)}.
 \label{a-finite-haar-analysis}
\end{equation}
With our positive Fourier sign, the analysis row of character $q$
pairs with the complex-conjugate wavelet character. This only permutes
the nonzero character labels $q\leftrightarrow p-q$ and leaves the
detail projectors unchanged.

The zero child-frequency channel retains a uniform superposition over
the new digit. Retaining standard residue zero instead keeps one
address out of each $p$ addresses. Both spaces have the same dimension,
but they are distinct subspaces. Their compressed dynamics and Schur
complements consequently need not agree. Digit reversal and child
Fourier analysis specify the transformations needed to compare them.



\section{Which Hamiltonians are compatible with refinement?}
\label{a-finite-hamiltonian}

An isometric address refinement does not automatically preserve a
Hamiltonian. The exact condition is a block condition on that
Hamiltonian. It can be checked without an approximation, a scaling
limit, or a physical interpretation of the new digit.

\begin{theorem}[Finite refinement classification]
\label{a-finite-hamiltonian-classification}
Let $V:\mathcal H_c\to\mathcal H_f$ be an isometry between finite
Hilbert spaces. Set $P=VV^*$ and $Q=I-P$. For self-adjoint operators
$H_c,H_f$, the following conditions are equivalent:
\begin{align}
 H_fV&=VH_c,\label{a-finite-hamiltonian-intertwine}\\
 QH_fP&=0,\qquad PH_fP=VH_cV^*.
 \label{a-finite-hamiltonian-blocks}
\end{align}
They hold exactly when, on $\Ran P\oplus\Ran Q$,
\begin{equation}
 H_f=(VH_cV^*)|_{\Ran P}\ \oplus\ H_Q,
 \qquad H_Q=(QH_fQ)|_{\Ran Q}.
 \label{a-finite-hamiltonian-direct-sum}
\end{equation}
The detail operator $H_Q$ may be any self-adjoint operator on
$\Ran Q$.
\end{theorem}

\begin{proof}
Right multiplication of \eqref{a-finite-hamiltonian-intertwine} by
$V^*$ gives $H_fP=VH_cV^*$. Left multiplication by $Q$ and $P$
gives \eqref{a-finite-hamiltonian-blocks}. Conversely $PV=V$ and
$V^*V=I$ imply
\[
 H_fV=(PH_fP+QH_fP)V=VH_cV^*V=VH_c.
\]
Self-adjointness gives $PH_fQ=(QH_fP)^*=0$. Thus both off-diagonal
blocks vanish, and the remaining blocks are exactly those in
\eqref{a-finite-hamiltonian-direct-sum}. Conversely that direct sum
has the two block properties. There is no further condition on the
self-adjoint detail block.
\end{proof}

The distinction between invariance and compression is essential.
Defining $H_c=V^*H_fV$ always matches the coarse block, but leaves
$QH_fV$ unconstrained. Thus compression alone does not establish
\eqref{a-finite-hamiltonian-intertwine}. When it does hold, the coarse
subspace is reducing, and fine evolution introduces no new detail
amplitude from a coarse initial state.

\begin{corollary}[Functional-calculus intertwining]
\label{a-finite-functional-calculus}
Under Theorem~\ref{a-finite-hamiltonian-classification}, for every
complex function $f$ defined on the finite union of the spectra,
\[
 f(H_f)V=Vf(H_c).
\]
In particular, for real $t$, real $\beta$, and
$z\notin\sigma(H_f)$,
\begin{align}
 e^{-\mathrm{i}tH_f}V&=Ve^{-\mathrm{i}tH_c},\nonumber\\
 e^{-\beta H_f}V&=Ve^{-\beta H_c},\nonumber\\
 (z-H_f)^{-1}V&=V(z-H_c)^{-1}.
 \label{a-finite-evolution-intertwining}
\end{align}
\end{corollary}

\begin{proof}
On $\Ran P$, $V$ is a unitary identification of $H_c$ with the
first block of $H_f$. The finite spectral decomposition therefore
maps each coarse spectral subspace to the corresponding first-block
spectral subspace. Applying $f$ to the eigenvalues gives the identity.
The three displayed functions are defined on the relevant spectra.
\end{proof}

The Euclidean operator $e^{-\beta H}$ in this corollary is a spectral
multiplier. It is not generally unitary, and a normalized state after
Euclidean evolution requires its own normalization. The identity does
not compare normalized Gibbs states of the entire fine and coarse
systems, since the fine partition function also counts the detail
block. For a compatible finite tower, the coarse propagator and its
embedded fine propagator agree exactly. This does not make an arbitrary
interacting fine Hamiltonian compatible.

\subsection{Hierarchical and translation-invariant examples}

Choose a real sequence of level energies $\epsilon_{p,r}$ and define
\begin{equation}
 H_{p,h}^{\mathrm{Haar}}=\sum_{r=0}^h
          \epsilon_{p,r}D_r^{(h)}.
 \label{a-finite-hierarchical-hamiltonian}
\end{equation}
Every wavelet in $\Ran D_r^{(h)}$ has energy $\epsilon_{p,r}$.
The projectors in \eqref{a-finite-detail-intertwining} show directly
that $H_{p,h+1}^{\mathrm{Haar}}V_{p,h}
=V_{p,h}H_{p,h}^{\mathrm{Haar}}$. The new detail energy
$\epsilon_{p,h+1}$ does not alter the embedded older levels.
Moreover the recursive Haar transform diagonalizes the level blocks:
\[
 \mathcal T_{p,h}H_{p,h}^{\mathrm{Haar}}\mathcal T_{p,h}^*
 =\sum_{r=0}^h\epsilon_{p,r}\Pi_r.
\]
This is exact for each finite $h$ and every finite real choice of the
energies. Their growth becomes a domain question only when an
infinite-depth operator is constructed.

Let the cyclic translation be $(S_{p,h}f)(x)=f(x-1)$, with
subtraction modulo $p^h$. Since a fine replicated function depends only
on the residue modulo $p^h$,
\[
 S_{p,h+1}V_{p,h}=V_{p,h}S_{p,h}.
\]
The adjoint translation obeys the same identity. Therefore the cyclic
Laplacian $L_{p,h}=2I-S_{p,h}-S_{p,h}^*$ intertwines, and for real
$\kappa_p$ the family
\begin{equation}
 H_{p,h}=H_{p,h}^{\mathrm{Haar}}+\kappa_p L_{p,h}
 \label{a-finite-compatible-laplacian}
\end{equation}
is self-adjoint and refinement compatible. Translation permutes the
level-$r$ cylinders, so it commutes with each $E_r^{(h)}$ and
$D_r^{(h)}$. The Laplacian thus preserves every detail subspace,
although a chosen spatially localized Haar wavelet need not be its
eigenvector.

This cyclic translation uses fixed unit displacement in the address
group. Multiplying it by a depth-dependent factor intended as a
physical inverse lattice spacing changes the tower. Such a factor is
not covered by \eqref{a-finite-compatible-laplacian} unless the
resulting operators still satisfy the stated intertwining equation.
Exact adic compatibility and continuum spatial consistency are
different requirements.

\subsection{A valuation-shell propagator}

For the pure hierarchical Hamiltonian, the finite spectral theorem gives
\begin{equation}
 U_{p,h}(t)=e^{-\mathrm{i}tH_{p,h}^{\mathrm{Haar}}}
 =\sum_{r=0}^he^{-\mathrm{i}t\epsilon_{p,r}}D_r^{(h)}.
 \label{a-finite-shell-propagator}
\end{equation}
The cylinder kernel is
\[
 E_r^{(h)}(x,y)=p^{-(h-r)}
              \mathbf1_{\{x\equiv y\, (\mathrm{mod}\,p^r)\}}.
\]
Consequently
\begin{align}
 U_{p,h}(t;x,y)
 &=p^{-h}e^{-\mathrm{i}t\epsilon_{p,0}}\nonumber\\
 &\quad+\sum_{r=1}^he^{-\mathrm{i}t\epsilon_{p,r}}
 \left[p^{-(h-r)}\mathbf1_{x\equiv y\,(\mathrm{mod}\,p^r)}
       -p^{-(h-r+1)}\mathbf1_{x\equiv y\,(\mathrm{mod}\,p^{r-1})}
 \right].
 \label{a-finite-shell-kernel}
\end{align}
Every indicator is determined by the truncated valuation
$\min\{v_p(x-y),h\}$, with the diagonal assigned value $h$.
The transition amplitude, and therefore the individual-address Born
probability, is constant on each such valuation shell. At a fixed
initial address $y$, the shell of valuation $k<h$ contains
$(p-1)p^{h-k-1}$ final addresses. Its total probability includes
this multiplicity. The diagonal has one address. Uniform probability
per address on a shell does not mean equal total probabilities of
different shells.



\subsection{Leakage detects failure of compatibility}
\label{a-finite-schur}

For a self-adjoint fine Hamiltonian with coarse projector $P$, set
$B=QH_fP$. If a unit vector $\psi$ lies in $\Ran P$, its initial
detail amplitude is
\[
 Qe^{-\mathrm{i}tH_f}\psi=-\mathrm{i}tB\psi+O(t^2),
 \qquad
 \norm{Qe^{-\mathrm{i}tH_f}\psi}^2
 =t^2\norm{B\psi}^2+O(|t|^3).
\]
These are finite Taylor expansions for fixed $H_f$ and $\psi$.
Thus a nonzero $B$ supplies a state with nonzero short-time leakage.
The two-dimensional example
\[
 H_f=\begin{pmatrix}0&\overline\gamma\\\gamma&0\end{pmatrix},
 \qquad P=\begin{pmatrix}1&0\\0&0\end{pmatrix}
\]
has compressed Hamiltonian zero but detail probability
$\sin^2(|\gamma|t)$ from the first basis state. It fails exact
refinement compatibility whenever $\gamma\neq0$.

For $z$ outside the spectrum of $H_Q=(QH_fQ)|_{\Ran Q}$, the
coarse Schur self-energy is
\begin{equation}
 \Sigma_P(z)=PH_fQ(z-H_Q)^{-1}QH_fP.
 \label{a-finite-self-energy}
\end{equation}
Compatibility makes it identically zero by setting $B=0$. This zero
is an algebraic statement, not a uniform bound on a resolvent near
its spectrum. A nonzero coupling requires the energy-dependent Schur
term when the detail sector is eliminated. The spectral part of this
monograph develops that elimination and its time-domain memory.

\adicfigure{05_projector_tower}{A finite compatible Hamiltonian has
zero blocks between its embedded coarse space and its new detail
space. In the plotted example, adding an off-diagonal coupling creates
detail probability from a coarse initial state. This finite example
illustrates the block criterion and supplies a negative control for
unconditional refinement claims.}

\section{Normalization and collective coupling}
\label{a-finite-normalization}

Refinement fixes the norm of an embedded packet. It does not fix an
interaction matrix element. The distinction becomes visible when one
coarse packet couples to many descendants rather than to one already
normalized aggregate packet. In this section $\hbar=1$, so the
Hamiltonian has units of inverse time.

\subsection{Descendant weights are fixed by the isometry}

After $w$ refinements a packet has $n=p^w$ orthonormal descendants.
Its normalized radial vector is
\begin{equation}
 |R_{p,w}\rangle=n^{-1/2}\sum_{\mathbf d}|\mathbf d\rangle.
 \label{a-finite-radial-vector}
\end{equation}
The norm is $n^{-1}\sum_{\mathbf d}1=1$, so each descendant has
amplitude modulus $p^{-w/2}$ and Born probability $p^{-w}$. This
reproduces \eqref{a-finite-nested-refinement}. The words ``radial'' and
``angular'' here refer to the uniform vector and its orthogonal
complement in a finite descendant space. They do not introduce a
spacetime radius or a rotation group.

Introduce a visible address $|c\rangle$ orthogonal to all descendants.
Suppose every descendant has the same coupling from $|c\rangle$,
with complex matrix element $\gamma_{p,w}$. The Hermitian interaction
is
\begin{equation}
 H_{\mathrm{int}}=
 \sum_{\mathbf d}
 \left(\gamma_{p,w}|\mathbf d\rangle\langle c|
       +\overline{\gamma_{p,w}}|c\rangle\langle\mathbf d|\right).
 \label{a-finite-leaf-interaction}
\end{equation}
The conjugate reverse coefficient is necessary when
$\gamma_{p,w}$ is not real. Define the collective coefficient
\begin{equation}
 G_{p,w}=\sqrt n\,\gamma_{p,w}.
 \label{a-finite-collective-coefficient}
\end{equation}
Then $H_{\mathrm{int}}|c\rangle=G_{p,w}|R_{p,w}\rangle$ and
$H_{\mathrm{int}}|R_{p,w}\rangle=
\overline{G_{p,w}}|c\rangle$. If
$|A\rangle=\sum_{\mathbf d}a_{\mathbf d}|\mathbf d\rangle$
is orthogonal to $|R_{p,w}\rangle$, its coefficients satisfy
$\sum_{\mathbf d}a_{\mathbf d}=0$. Hence
$H_{\mathrm{int}}|A\rangle=0$. Only the radial vector is bright
under this equal-leaf interaction.

This statement concerns the interaction block. A separate descendant
Hamiltonian can mix radial and angular vectors. To obtain an invariant
two-dimensional radial subsystem, the full descendant Hamiltonian
must preserve the radial line, and the full Hamiltonian must leave its
angular complement uncoupled from $|c\rangle$. A degenerate hidden
energy $\epsilon I$ is one sufficient choice.

\subsection{The exact bright two-level problem}

Assume those invariance conditions and order the basis as
$(|c\rangle,|R_{p,w}\rangle)$. With visible energy $E_c$ and
radial energy $\epsilon$, the exact bright block is
\begin{equation}
 H_{\mathrm{bright}}=
 \begin{pmatrix}E_c&\overline G_{p,w}\\G_{p,w}&\epsilon\end{pmatrix}.
 \label{a-finite-bright-block}
\end{equation}
Set $\overline E=(E_c+\epsilon)/2$,
$\delta=(E_c-\epsilon)/2$, and
$\Omega=(\delta^2+|G_{p,w}|^2)^{1/2}$. Then
$(H_{\mathrm{bright}}-\overline E I)^2=\Omega^2I$, giving
\begin{equation}
 e^{-\mathrm{i}tH_{\mathrm{bright}}}
 =e^{-\mathrm{i}t\overline E}
 \left[I\cos(\Omega t)
 -\mathrm{i}\frac{H_{\mathrm{bright}}-\overline E I}{\Omega}
              \sin(\Omega t)\right].
 \label{a-finite-bright-propagator}
\end{equation}
At $\Omega=0$ this formula is understood by continuous extension.
Starting in $|c\rangle$, the radial probability is therefore
\begin{equation}
 \Pr(R_{p,w};t)=
 \frac{|G_{p,w}|^2}{\delta^2+|G_{p,w}|^2}
       \sin^2\!\left(t\sqrt{\delta^2+|G_{p,w}|^2}\right).
 \label{a-finite-radial-probability}
\end{equation}
The probability is defined to be zero when $\Omega=0$, equivalently
when both $\delta$ and $G_{p,w}$ vanish. On resonance $E_c=\epsilon$ it reduces to
$\sin^2(|G_{p,w}|t)$. The individual descendant probabilities are
$\Pr(R_{p,w};t)/p^w$ when the hidden amplitude stays radial.

Eliminating the hidden radial level gives
\begin{equation}
 \Sigma_{p,w}(z)=\frac{|G_{p,w}|^2}{z-\epsilon},
 \qquad z\neq\epsilon.
 \label{a-finite-radial-self-energy}
\end{equation}
This is the exact finite Schur term. The full visible Green function,
when its denominator is nonzero, is
$[z-E_c-\Sigma_{p,w}(z)]^{-1}$. It is distinct from the self-energy
itself, and the locations of its poles depend on both sectors.

\subsection{The coupling exponent is a model choice}

Let $g\in\C$ be fixed and choose
$\gamma_{p,w}=g p^{-\alpha w}$ for a real exponent $\alpha$.
Then
\begin{equation}
 G_{p,w}=g p^{(1/2-\alpha)w},\qquad
 |G_{p,w}|^2=|g|^2p^{(1-2\alpha)w}.
 \label{a-finite-coupling-exponent}
\end{equation}
The exact branching multiplicity produces the factor $p^{w/2}$.
Thus the collective coefficient grows with $w$ for $\alpha<1/2$,
is independent of $w$ for $\alpha=1/2$, and decays for
$\alpha>1/2$. The same cases govern the operator norm of this
equal-leaf interaction: its two radial eigenvalues are
$\pm|G_{p,w}|$, so $\norm{H_{\mathrm{int}}}=|G_{p,w}|$.
They do not determine the boundedness of a different infinite operator
with many interacting depth sectors.

\begin{center}
\begin{tabular}{@{}lccc@{}}
\toprule
Per-leaf choice & $\alpha$ & $|G_{p,w}|$ & $|G_{p,w}|^2$\\
\midrule
Constant matrix element & $0$ & $|g|p^{w/2}$ & $|g|^2p^w$\\
Constant collective coupling & $1/2$ & $|g|$ & $|g|^2$\\
Uniform measure coefficient & $1$ & $|g|p^{-w/2}$ & $|g|^2p^{-w}$\\
\bottomrule
\end{tabular}
\end{center}

The measure coefficient $\alpha=1$ assigns the number $p^{-w}$ to
each per-leaf matrix element after factoring out $g$. It gives a
collective amplitude of size $p^{-w/2}$ and a self-energy numerator of
size $p^{-w}$. The self-energy estimate is uniform only where
$|z-\epsilon|$ stays bounded below. Likewise
\[
 \Pr(R_{p,w};t)\leq |g|^2t^2p^{-w}
 \quad\text{when }\alpha=1
\]
follows from \eqref{a-finite-radial-probability} and
$|\sin u|\leq|u|$. It is a fixed-time bound. On resonance a time
of order $p^{w/2}/|g|$ can still transfer order-one probability.

A model that directly couples $|c\rangle$ to one normalized hidden
packet with matrix element $g p^{-w}$ instead has a collective
coefficient of that same size and a self-energy numerator
$|g|^2p^{-2w}$. It is a valid two-level model, but it is not the
equal-leaf model \eqref{a-finite-leaf-interaction}. The number of
descendant addresses alone does not select between these interaction
geometries. A physical refinement law must specify the coupling,
the retained hidden modes, their energies, and the observation scale.


